\documentclass[11pt]{article}
\usepackage{amsmath,hyperref}
\title{An affordable send can consume the budget of an older live duty}
\date{}
\begin{document}\maketitle
\paragraph{Scope.}
A first old report actually arrives, but a second old report remains due.
In the pinned native virtual-testbed broker, a new report can be affordable
now and still spend the budget that the older follow-up will need. Joint
admission uses a qualified clock and the admitted scope of future old
effects. This is source-level civil simulation, not measured network
behavior or an independent case.

\paragraph{Mechanism.}
Use the selected
\href{https://github.com/osrf/subt/blob/083c6a8acb0c25aa1a35e34517e628cd3b52ae77/subt-communication/subt_communication_broker/src/subt_communication_broker.cpp}
{native broker dispatch} and
\href{https://github.com/osrf/subt/blob/083c6a8acb0c25aa1a35e34517e628cd3b52ae77/subt-communication/subt_communication_model/src/subt_communication_model.cpp}
{native radio calculation}.
The first old report of $a$ bytes is sent at time zero, the second of
$b$ at known $h<1$. A new report of $x$ bytes uses the first report's
sender and receiver. Its issue time is chosen from declared candidates;
an old attempt comes first at a tie. The second old report can share
the first report's physical sender, recipient, both or neither.
Actual receiver callbacks count as payments; independent work has a
separate pair of endpoints and continues. Paths in this fixture are
reliable except for budget refusal. The checker must choose a single
time valid for all admitted bindings, not select using a hidden
later source identity. A stale disjointness assertion has no
exclusion power.

\paragraph{Uniform-rate calculation.}
Assume $a+b\le C$, so the old duties are feasible without a new one.
With an early new send before $h$ and $a+x\le C<a+b+x$, new payment
can make the older second send fail on any shared resource. Reversing
the order reverses which report misses, not whether they fit jointly.
A genuinely disjoint tail avoids this failure. For $h\le t\le1$, a
shared budget still retains the first old charge; for $1<t\le h+1$
only the second charge remains and $b+x\le C$ suffices. If that bound
fails, a shared budget requires waiting beyond $h+1$. The native
window includes a charge at equality.
For $(C,a,b,x,h)=(5,2,2,3,0.75)$, known disjointness permits a new
send at $0.5$. Unknown overlap fails that deadline but works at the
first tested floating-point value strictly above one. In continuous
time one is an infimum, not an acceptable earliest timestamp.

\paragraph{Source-specific capacity and failed transport.}
At equal sending rates the three overlapping modes have the same joint
verdict. At unequal rates they do not. The initial hypothesis that TX
and RX separately used their own node's configured capacity failed
the native replay. Source review and controlled substitution identified
the reason: the broker passes the sending node's configuration into
both budget checks, while separately retaining the TX and RX histories
on the actual nodes. A receiver ledger's effective limit can therefore
change with the sender of the current report.
For $a=b=x=2$ and sending rates $(P,Q)=(5,6)$, early new payment
prevents the old follow-up in the sender-only overlap but not the
receiver-only overlap. With sending rates $(6,5)$ the verdict reverses.
Changing the receive-only node's own configured rate from 3 to 7
does not change this test. The uniform overlap flag alone thus cannot
justify transport to unequal configured rates. No broader claim about
the larger simulator's receiver configuration follows.

\paragraph{Verification and residual.}
A separate time-indexed ledger simulator agrees with 43,680 uniform
and 240 unequal native effect rows, and the exact competing cut
enumeration agrees on 87,360 selections. Out of the raw uniform cases,
3,150 early candidate sends are currently affordable but cause a
required old receipt to fail. Lean checks three arithmetic statements.
This is a different source mechanism from a standing warehouse shelf:
the debit can truly expire without changing scope. The exact direct
control still ties. A real transfer owes an attributable scope producer,
visibility/clock qualification and access for the acting recipients;
this is not a contract acceptance or a fault disposition.
\end{document}
